%% Small-base presheaves, pointwise constructions, and subobjects.

\subsection{Presheaves and their pointwise constructions}
\label{sec:p10-presheaf-categories}

\begin{hbdefinition}{def:p10-presheaf-category}{The category of presheaves on a small category}
Fix a category \(\mathcal C\) whose object and arrow sets belong
to the fixed universe.  A presheaf on \(\mathcal C\) is a functor
\[
                         X:\mathcal C^{\mathrm{op}}\to\mathsf{Set}.
\]
Thus \(u:d\to c\) induces a restriction function
\(X(u):X(c)\to X(d)\).
For \(v:e\to d\), the functor laws are
\[
                    X(1_c)=1_{X(c)},\qquad
                    X(u v)=X(v)X(u).
\]
A map of presheaves \(\alpha:X\to Y\) is a natural
transformation: its functions \(\alpha_c:X(c)\to Y(c)\) satisfy
\[
                         Y(u)\alpha_c=\alpha_d X(u).
\]
Write
\(\widehat{\mathcal C}=[\mathcal C^{\mathrm{op}},\mathsf{Set}]\)
for their category, with componentwise identities and composition.
The category itself may be large in the fixed universe.

For \(c\in\mathcal C\), the representable presheaf is
\(\mathrm y c=\mathcal C(-,c)\).
An arrow \(u:d\to c\) induces
\[
 \mathrm y(u):\mathrm y d\longrightarrow\mathrm y c,\qquad
 (\mathrm y(u))_e(v)=u v\quad(v:e\to d).
\]
This is the covariant Yoneda assignment on representing objects;
each \(\mathrm y c\) is contravariant in its argument.
\index{presheaf}\index{presheaf category}\index{Yoneda embedding}
\end{hbdefinition}

\begin{hbdefinition}{def:p10-subpresheaf}{Subpresheaves and their order}
A subpresheaf \(S\) of \(X\) is a family of subsets
\(S(c)\subseteq X(c)\) such that
\[
                         X(u)[S(c)]\subseteq S(d)
                         \quad(u:d\to c).
\]
Its restriction maps are the restrictions of those of \(X\).
Write \(\operatorname{Sub}(X)\) for the set of such families,
ordered by componentwise inclusion.
This is a small set: it is a subset of
\[
                       \prod_{c\in\operatorname{Ob}\mathcal C}
                                      \mathcal P(X(c)).
\]
The compatibility conditions select a subset of this product.
The theorem below identifies these families with representatives
of categorical subobjects of \(X\).
\index{subpresheaf}\index{subobject!of a presheaf}
\end{hbdefinition}

\begin{hbtheorem}{thm:p10-presheaf-limits}{proved=annex:p10.psh.1}{Small limits and colimits of presheaves}
The category \(\widehat{\mathcal C}\) is locally small.
It has all small limits and colimits, computed pointwise.
More precisely, for a small diagram
\(D:J\to\widehat{\mathcal C}\), the component at \(c\) of
its limiting cone or colimiting cocone is the corresponding
universal cone or cocone for the set diagram \(j\mapsto D(j)(c)\).
Canonical set constructions supply the restriction maps without
choosing an unrelated limit or colimit at every object.

In particular, the terminal presheaf \(1\) has value \(\{*\}\)
at every object, and the initial presheaf \(0\) has value
\(\varnothing\).
Products have
\((X\times Y)(c)=X(c)\times Y(c)\), and the equalizer of
\(\alpha,\beta:X\to Y\) has component
\[
                         \{x\in X(c):\alpha_c(x)=\beta_c(x)\}.
\]
Coproducts and coequalizers are the corresponding pointwise
disjoint unions and quotients.
If \(\mathcal C\) is empty, its presheaf category has one object
and one arrow; the same constructions give its initial and
terminal object.
\end{hbtheorem}

\begin{hbtheorem}{thm:p10-presheaf-mono-epi}{proved=annex:p10.psh.2}{Maps and subobjects of presheaves}
For a natural map \(m:X\to Y\):
\[
 \begin{aligned}
 m\text{ is monic}
      &\quad\Longleftrightarrow\quad
                    m_c\text{ is injective for every }c,\\
 m\text{ is epic}
      &\quad\Longleftrightarrow\quad
                    m_c\text{ is surjective for every }c,\\
 m\text{ is an isomorphism}
      &\quad\Longleftrightarrow\quad
                    m_c\text{ is bijective for every }c.
 \end{aligned}
\]
Its componentwise images form a subpresheaf
\[
                    \operatorname{Im}(m)(c)=m_c[X(c)]
\]
and give a factorization
\[
                     X\xrightarrow{q}\operatorname{Im}(m)
                                      \xrightarrow{i}Y,
\]
where \(q\) is epic and \(i\) is the inclusion and is monic.
If \(m\) is monic, \(q\) is an isomorphism over \(Y\).

Two monomorphisms into \(Y\) are isomorphic over \(Y\) exactly
when their image subpresheaves agree.  One such monomorphism
factors through another exactly when its image is contained
in the other's image.
Thus \(\operatorname{Sub}(Y)\) gives a canonical small set of
representatives for the ordered categorical subobjects.
A map that is both monic and epic is an isomorphism.
No assertion that every epimorphism has a natural section is made.
\end{hbtheorem}

The pointwise constructions and componentwise map criteria are
stated in \cite[\S1, p.~3]{MoerdijkVanOosten2007Topos}.
The annex supplies the complete constructions and cancellation
proofs, including the converse epimorphism argument through a
pointwise pushout.
For the empty base there are no components; the component
criteria are vacuous and hold for the unique arrow.
